% P5 of the public test corpus (docs/DESIGN.md#testing). Written by the project; dedicated to the public domain
% (CC0 1.0). The look of a Nature or npj manuscript on the standard article class (no publisher class):
% unnumbered headings, superscript numbered citations, line numbers, double spacing, a Methods section and
% figure legends at the end.
\documentclass[11pt]{article}
\usepackage[T1]{fontenc}
\usepackage[utf8]{inputenc}
\usepackage[a4paper,margin=2.5cm]{geometry}
\usepackage{amsmath}
\usepackage{graphicx}
\usepackage{authblk}
\usepackage{lineno}
\usepackage{setspace}
\usepackage[super,sort&compress,comma]{natbib}
\usepackage[hidelinks]{hyperref}
\usepackage{titlesec}
\titleformat{\section}{\normalfont\large\bfseries}{}{0pt}{}

\title{Phonon lifetimes from a two-temperature model}
\author[1]{First Author}
\author[1,2]{Second Author}
\affil[1]{Department of Physics, Example University}
\affil[2]{Institute for Test Documents}
\date{}

\linenumbers
\doublespacing

\begin{document}
\maketitle

\begin{abstract}
Short abstracts of about 150 words state the problem, the approach and the main result. Here a
two-temperature model with one coupling constant reproduces measured phonon lifetimes across three
materials, with no other fitted parameter.
\end{abstract}

\section{Introduction}
Electron--phonon coupling sets the rate at which a hot electron gas heats the lattice\cite{allen1987}. A
two-temperature model describes the exchange with one coupling constant $G$:
\begin{equation}
  C_e \frac{\mathrm{d}T_e}{\mathrm{d}t} = -G\,(T_e - T_l), \qquad
  C_l \frac{\mathrm{d}T_l}{\mathrm{d}t} = G\,(T_e - T_l).
\end{equation}

\section{Results}
The model fits all three materials with a single value of $G$ (Fig.~1)\cite{allen1987,ziman1960}.

\section{Discussion}
The agreement suggests that one coupling constant suffices in this temperature range\cite{ziman1960}.

\section{Methods}
\paragraph{Model.} The coupled equations were integrated with a fourth-order Runge--Kutta method.
\paragraph{Data availability.} All data are synthetic and generated by the code that accompanies the text.

\bibliographystyle{unsrtnat}
\bibliography{manuscript}

\section*{Figure legends}
\textbf{Figure 1 \textbar\ Lattice temperature against time.} The model (lines) against the synthetic data
(points) for three materials.

\end{document}
